% 原创教学示意；共享消息函数与求和聚合，不依赖邻居枚举次序。
\begin{tikzpicture}[x={\dimexpr\linewidth/169\relax},y=1mm,
  font=\figurefont\footnotesize,text=NNInk]
  \node[nn hidden] (i) at (24,0) {$i$};
  \node[nn input] (j) at (5,20) {$j$};
  \node[nn input] (k) at (5,-20) {$k$};
  \node[nn input] (r) at (40,-20) {$r$};
  \foreach \v in {j,k,r}{\draw[nn flow] (\v)--(i);}
  \draw[nn link] (47,-29)--(47,32);
  \foreach \v/\yy in {j/18,k/0,r/-18}{
    \node[anchor=east] (h\v) at (60,\yy) {$\bm h_{\v}^{(l)}$};
    \node[nn operator,minimum width=16mm] (m\v) at (75,\yy) {$M^{(l)}$};
    \draw[nn flow] (h\v)--(m\v);
  }
  \node[nn pointwise,minimum size=8mm] (sum) at (107,0) {$\sum$};
  \foreach \v in {j,k,r}{\draw[nn operator path] (m\v.east)--(sum.west);}
  \node[nn operator,minimum width=18mm] (update) at (135,0) {$U^{(l)}$};
  \node[anchor=south] (self) at (135,25) {$\bm h_i^{(l)}$};
  \draw[nn operator path] (self)--(update.north);
  \draw[nn operator path] (sum)--(update);
  \node[anchor=west] (out) at (153,0) {$\bm h_i^{(l+1)}$};
  \draw[nn flow,draw=NNOutput] (update)--(out);
\end{tikzpicture}
